\chapter{Hamilton-Jacobi Theory And Adiabatic Invariants}
\label{chap:hamilton-jacobi}

\section{Why Hamilton-Jacobi Belongs Here}
\label{sec:why-hamilton-jacobi}

Hamilton-Jacobi theory belongs naturally between canonical transformations and
integrable systems. The Hamilton-Jacobi equation is not merely another way to
solve special examples. It explains why a canonical transformation generated by
one scalar function can turn a Hamiltonian flow into trivial motion, and it
gives the conceptual bridge from Hamiltonian mechanics to action-angle
variables.

\begin{table}[htbp]
\centering
\caption{Notation for Hamilton-Jacobi theory and adiabatic invariants.}
\label{tab:hj-notation}
\small
\begin{tabular}{p{0.18\textwidth}p{0.48\textwidth}p{0.25\textwidth}}
\toprule
Symbol & Meaning & Type\\
\midrule
\(q^i,p_i\) & old canonical coordinates and momenta & local coordinates on \(T^*Q\)\\
\(Q^i,P_i\) & new canonical coordinates and momenta & transformed phase-space coordinates\\
\(S(q,\alpha,t)\) & Hamilton's principal function or complete integral & generating function\\
\(\alpha_i\) & constants used as new momenta \(P_i\) & integration constants\\
\(\beta^i\) & constants \(\partial S/\partial\alpha_i\) & new coordinates\\
\(W(q,\alpha)\) & Hamilton's characteristic function & time-independent part of \(S\)\\
\(E\) & conserved energy in autonomous systems & scalar\\
\(I\) & one-dimensional action variable & \((2\pi)^{-1}\oint p\,\dd q\)\\
\(\theta\) & angle variable conjugate to \(I\) & \(2\pi\)-periodic\\
\(L(t)\) & slowly varying wall separation in the box example & length\\
\(\epsilon_{\mathrm{ad}}\) & adiabatic small parameter & period / driving time scale\\
\bottomrule
\end{tabular}
\end{table}

\paragraph{Roadmap.}
The chapter starts with generating functions because the Hamilton-Jacobi
equation is best understood as a canonical transformation problem. A complete
integral \(S(q,\alpha,t)\) is not just a solution of a PDE; it is a device that
turns the original Hamiltonian flow into constant new momenta and linearly
recoverable old coordinates. The autonomous case introduces the characteristic
function \(W\), the action integral connects this to invariant tori, and the
moving-wall example shows what remains true when parameters vary slowly rather
than exactly stay fixed.

\paragraph{Two meanings of ``action''.}
The word action appears in two related but distinct senses. Hamilton's
principle uses the time integral \(\int L\,\dd t\) as a functional on paths.
Action-angle theory uses \(I=(2\pi)^{-1}\oint p\,\dd q\) as a coordinate
labelling a closed orbit or invariant torus. Hamilton-Jacobi theory is the
bridge: its generating function accumulates the canonical one-form
\(p_i\,\dd q^i-H\,\dd t\), and its period integrals become the action
coordinates used in integrability and adiabatic invariance.

\section{Generating Functions And The Hamilton-Jacobi Equation}
\label{sec:hj-generating-function}

The key idea is to choose new canonical coordinates in which the motion is as
simple as possible. If the new Hamiltonian vanishes, then the new coordinates
and momenta are constants. The price is that the generating function must solve
a nonlinear first-order PDE. That PDE is the Hamilton-Jacobi equation.

Let \((q,p)\) be old canonical coordinates and \((Q,P)\) new canonical
coordinates. A time-dependent type-II generating function
\[
  S=S(q,P,t)
\]
defines a canonical transformation by
\[
  p_i=\frac{\partial S}{\partial q^i},
  \qquad
  Q^i=\frac{\partial S}{\partial P_i}.
\]
These formulas are not arbitrary. They are chosen so that
\[
  \dd S=p_i\,\dd q^i+Q^i\,\dd P_i+\frac{\partial S}{\partial t}\,\dd t.
\]
Since \(Q^i\dd P_i=\dd(P_iQ^i)-P_i\dd Q^i\), this can be rearranged as
\[
  p_i\,\dd q^i
  =
  P_i\,\dd Q^i
  +
  \dd(S-P_iQ^i)
  -
  \frac{\partial S}{\partial t}\,\dd t.
\]
The phase-space action transforms as
\[
  p_i\,\dd q^i-H\,\dd t
  =
  P_i\,\dd Q^i-K\,\dd t+\dd(S-P_iQ^i),
\]
where \(K(Q,P,t)\) is the new Hamiltonian. Comparing coefficients gives
\[
  K(Q,P,t)
  =
  H\!\left(q,\frac{\partial S}{\partial q},t\right)
  +
  \frac{\partial S}{\partial t}.
\]
The Hamilton-Jacobi strategy is to choose \(S\) so that the new Hamiltonian is
as simple as possible. The strongest choice is
\[
  K=0.
\]
Then \(P_i\) and \(Q^i\) are all constants along the new motion. Writing
\(P_i=\alpha_i\), the required equation is
\[
  \boxed{
  \frac{\partial S}{\partial t}
  +
  H\!\left(q,\frac{\partial S}{\partial q},t\right)=0.
  }
\]
This is the Hamilton-Jacobi equation.

\begin{figure}[htbp]
\centering
\begin{tikzpicture}[scale=1.0, >=Latex]
  \draw[->] (-0.2,0) -- (6.4,0) node[right] {\(q^1\)};
  \draw[->] (0,-0.2) -- (0,3.6) node[above] {\(q^2\)};
  \draw[blue!70!black, thick]
    (0.6,0.6) .. controls (1.8,1.1) and (3.1,1.0) .. (5.4,0.65);
  \draw[blue!70!black, thick]
    (0.8,1.35) .. controls (2.0,1.85) and (3.35,1.7) .. (5.65,1.35);
  \draw[blue!70!black, thick]
    (1.0,2.1) .. controls (2.25,2.65) and (3.65,2.45) .. (5.9,2.1);
  \draw[red!75!black, thick, ->]
    (1.45,0.78) .. controls (1.65,1.28) and (1.9,1.68) .. (2.2,2.24);
  \draw[red!75!black, thick, ->]
    (3.0,0.9) .. controls (3.15,1.35) and (3.35,1.75) .. (3.72,2.42);
  \draw[red!75!black, thick, ->]
    (4.65,0.76) .. controls (4.85,1.20) and (5.05,1.55) .. (5.38,2.15);
  \draw[->, black!70] (3.35,1.7) -- (3.85,2.05)
    node[right] {\(p=\nabla_qS\)};
  \node[blue!70!black, align=center] at (1.05,3.15)
    {level sets\\\(S(q,t)=\text{constant}\)};
  \node[red!75!black, align=center] at (5.15,3.05)
    {characteristics\\classical paths};
\end{tikzpicture}
\caption{Hamilton-Jacobi geometry. At a fixed time the level sets of
\(S(q,t)\) are wavefronts in configuration space; the momentum covector
\(p=\nabla_qS\) is normal to those wavefronts, and Hamilton's equations carry
these normals along the characteristic curves.}
\label{fig:hj-wavefront-characteristics}
\end{figure}

\begin{definition}[Complete integral]
A function \(S(q,\alpha,t)\) is a complete integral of the Hamilton-Jacobi
equation for an \(n\)-degree-of-freedom system if it depends on \(n\) independent
parameters \(\alpha_1,\ldots,\alpha_n\) and satisfies
\[
  \det\left(
    \frac{\partial^2S}{\partial q^i\partial\alpha_j}
  \right)\neq0
\]
on the region under consideration. The nondegeneracy condition is exactly what
lets \((q,\alpha)\mapsto(q,p)\) define a local canonical transformation.
\end{definition}

Once a complete integral is known, the solution is obtained algebraically:
\[
  p_i(t)=\frac{\partial S}{\partial q^i}(q(t),\alpha,t),
  \qquad
  \beta^i=\frac{\partial S}{\partial\alpha_i}(q(t),\alpha,t),
\]
where \(\alpha_i\) and \(\beta^i\) are constants fixed by initial data. This is
why Hamilton-Jacobi theory is sometimes called integration by quadrature: the
hard part is constructing \(S\); after that the motion is recovered by
differentiation and inversion.

The nondegeneracy condition in the definition is essential. Without it,
\(\alpha\) would not provide a valid set of new momenta and the equations
\(\partial S/\partial\alpha=\beta\) could fail to determine \(q(t)\). Thus a
Hamilton-Jacobi solution is useful not merely because it solves a PDE, but
because it generates a canonical coordinate system adapted to the flow.

\section{Autonomous Systems And The Characteristic Function}
\label{sec:hj-autonomous}

If \(H(q,p)\) is independent of time, seek a separated form
\[
  S(q,\alpha,t)=W(q,\alpha)-Et,
\]
where one of the constants \(\alpha_i\) is the energy \(E\). The
Hamilton-Jacobi equation becomes
\[
  \boxed{
  H\!\left(q,\frac{\partial W}{\partial q}\right)=E.
  }
\]
This is the time-independent Hamilton-Jacobi equation. The function \(W\) is
Hamilton's characteristic function.

For one degree of freedom with Hamiltonian
\[
  H(q,p)=\frac{p^2}{2m}+V(q),
\]
the equation reads
\[
  \frac1{2m}\left(\frac{\dd W}{\dd q}\right)^2+V(q)=E.
\]
Thus
\[
  \frac{\dd W}{\dd q}=p(q,E)=\pm\sqrt{2m(E-V(q))},
\]
and
\[
  W(q,E)=\int^q p(q',E)\,\dd q'.
\]
This formula is simple but profound: the phase of the Hamilton-Jacobi function
is the accumulated canonical one-form \(p\,\dd q\).

\section{Action As A Period Integral}
\label{sec:hj-action-period}

Suppose a one-degree-of-freedom autonomous system has a bounded orbit between
turning points \(q_-(E)\) and \(q_+(E)\). The action variable is
\[
  \boxed{
  I(E)=\frac{1}{2\pi}\oint p\,\dd q
  =
  \frac1{\pi}\int_{q_-(E)}^{q_+(E)}
  \sqrt{2m(E-V(q))}\,\dd q.
  }
\]
The factor \(1/\pi\) appears because the closed orbit has a positive-momentum
branch and a negative-momentum branch.

\begin{figure}[htbp]
\centering
\begin{tikzpicture}[scale=1.0, >=Latex]
  \draw[->] (-0.4,0) -- (6.6,0) node[right] {\(q\)};
  \draw[->] (0,-2.0) -- (0,2.2) node[above] {\(p\)};
  \draw[densely dotted] (1.0,-1.7) -- (1.0,1.8);
  \draw[densely dotted] (5.3,-1.7) -- (5.3,1.8);
  \node[below] at (1.0,0) {\(q_-\)};
  \node[below] at (5.3,0) {\(q_+\)};
  \draw[blue!70!black, thick, ->]
    (1.0,0) .. controls (2.0,1.55) and (4.2,1.55) .. (5.3,0);
  \draw[blue!70!black, thick, ->]
    (5.3,0) .. controls (4.2,-1.55) and (2.0,-1.55) .. (1.0,0);
  \node[blue!70!black] at (3.2,1.65) {\(+p(q,E)\)};
  \node[blue!70!black] at (3.2,-1.65) {\(-p(q,E)\)};
  \node[align=center] at (3.2,0.55)
    {enclosed area\\\(\displaystyle \oint p\,\dd q=2\pi I\)};
\end{tikzpicture}
\caption{For a one-dimensional bound orbit the action is the signed
phase-space area enclosed by the energy curve divided by \(2\pi\). The turning
points contribute no endpoint term when differentiating with respect to
energy because \(p(q_\pm,E)=0\).}
\label{fig:hj-action-loop}
\end{figure}

The period follows from differentiating \(I(E)\). The moving endpoint terms
vanish because the momentum is zero at the turning points:
\[
  p(q_\pm(E),E)=0.
\]
Thus
\[
  \frac{\dd I}{\dd E}
  =
  \frac1{\pi}\int_{q_-}^{q_+}
  \frac{m\,\dd q}{\sqrt{2m(E-V(q))}}
  =
  \frac1{2\pi}T(E).
\]
Therefore
\[
  \boxed{
  \omega(E)=\frac{\dd H}{\dd I}
  =
  \left(\frac{\dd I}{\dd E}\right)^{-1},
  \qquad
  T(E)=2\pi\frac{\dd I}{\dd E}.
  }
\]
This formula is the one-dimensional seed of action-angle theory.

\begin{example}[Harmonic oscillator]
For
\[
  H=\frac{p^2}{2m}+\frac12m\omega_0^2q^2,
\]
the phase curve is an ellipse:
\[
  \frac{p^2}{2mE}+\frac{q^2}{2E/(m\omega_0^2)}=1.
\]
The area enclosed by the ellipse is
\[
  \pi\sqrt{2mE}\sqrt{\frac{2E}{m\omega_0^2}}
  =
  \frac{2\pi E}{\omega_0}.
\]
Thus
\[
  I=\frac{1}{2\pi}\operatorname{Area}=\frac{E}{\omega_0},
  \qquad
  H(I)=\omega_0 I.
\]
The conjugate angle advances uniformly:
\[
  \dot\theta=\frac{\partial H}{\partial I}=\omega_0.
\]
\end{example}

\section{Hamilton-Jacobi And Wave Optics}
\label{sec:hj-wave-analogy}

Hamilton-Jacobi theory is also the classical side of a broader ray-wave
analogy. The analogy begins with the eikonal approximation. If a wave has
rapidly oscillating phase
\[
  \psi(q,t)\approx A(q,t)e^{\frac{i}{\hbar}S(q,t)},
\]
then substituting into the Schrodinger equation and keeping the leading power
of \(\hbar\) gives
\[
  \partial_tS+H\left(q,\partial_qS\right)=0.
\]
Thus classical mechanics is the short-wavelength or stationary-phase limit of
wave mechanics. In optics, rays are normals to wavefronts. In mechanics,
trajectories are characteristics of the Hamilton-Jacobi equation, and the
momentum is the phase gradient:
\[
  p_i=\partial_iS.
\]

\section{Adiabatic Invariant: Particle Between Slowly Moving Walls}
\label{sec:adiabatic-box}

A clean adiabatic-invariant model is a particle that bounces elastically
between a fixed wall and a slowly moving wall. Let \(L(t)\) be the wall
separation and \(p(t)=m v(t)>0\) the magnitude of the particle momentum between
collisions.
The example is intentionally elementary: no differential geometry is needed to
see the invariant. The point is that the same quantity obtained from collision
kinematics is also the action integral \((2\pi)^{-1}\oint p\,\dd q\). This is
why adiabatic invariance belongs with action-angle variables.

\begin{assumption}[Adiabatic wall motion]
The wall speed is small compared with the particle speed:
\[
  \epsilon_{\mathrm{ad}}
  =
  \frac{|\dot L|}{v}\ll1.
\]
Over one round trip, the wall speed can be treated as approximately constant.
\end{assumption}

First derive the invariant by elementary collision kinematics. If the moving
wall approaches the fixed wall with speed \(u=-\dot L>0\), an elastic collision
with that wall increases the particle speed by approximately \(2u\). During one
round trip,
\[
  \Delta t=\frac{2L}{v},
  \qquad
  \Delta L=-u\Delta t=-\frac{2Lu}{v}.
\]
The momentum change is
\[
  \Delta p=2mu=\frac{2pu}{v}.
\]
Eliminating \(u/v\) gives
\[
  \Delta p=-\frac{p}{L}\Delta L.
\]
Therefore
\[
  \Delta(pL)=L\Delta p+p\Delta L=0
\]
to leading adiabatic order. Hence
\[
  \boxed{p(t)L(t)\approx \text{constant}.}
\]

The action interpretation explains why this result is general. For a particle
in a one-dimensional box, the phase-space loop consists of a branch with
momentum \(+p\) from \(0\) to \(L\) and a branch with momentum \(-p\) from
\(L\) back to \(0\). Hence
\[
  \oint p\,\dd q=2pL,
  \qquad
  I=\frac{pL}{\pi}.
\]
Slow parameter changes preserve \(I\) to adiabatic accuracy, so \(pL\) is the
adiabatic invariant. The energy scales as
\[
  E(t)=\frac{p(t)^2}{2m}
  =
  E(0)\frac{L(0)^2}{L(t)^2}.
\]

\begin{figure}[htbp]
\centering
\begin{tikzpicture}[scale=1.0, >=Latex]
  \draw[->] (-0.4,0) -- (6.4,0) node[right] {\(q\)};
  \draw[->] (0,-1.9) -- (0,2.1) node[above] {\(p\)};
  \draw[thick] (0.8,-1.25) rectangle (5.4,1.25);
  \draw[blue!75!black, thick, ->] (0.8,1.25) -- (5.4,1.25);
  \draw[blue!75!black, thick, ->] (5.4,-1.25) -- (0.8,-1.25);
  \draw[dashed] (5.4,-1.55) -- (5.4,1.55);
  \draw[->, red!75!black, thick] (5.4,1.65) -- (4.8,1.65)
    node[left] {slow wall motion};
  \node[below] at (0.8,0) {\(0\)};
  \node[below] at (5.4,0) {\(L(t)\)};
  \node[right] at (5.4,1.25) {\(+p\)};
  \node[right] at (5.4,-1.25) {\(-p\)};
  \node[align=center] at (3.1,0)
    {area \(=2pL\)\\\(I=pL/\pi\)};
\end{tikzpicture}
\caption{The adiabatic box in phase space. Slow wall motion changes the
rectangle's width \(L\) and height \(p\), but preserves the enclosed area
\(2pL\) to leading adiabatic order.}
\label{fig:adiabatic-box-action}
\end{figure}

This example also explains the limitation of the word ``invariant.'' The
quantity \(I\) is not exactly conserved for arbitrary wall motion. It is
preserved to high accuracy when the parameter changes slowly compared with the
orbital period, and the accuracy improves as the ratio of time scales becomes
smaller. Adiabatic invariance is therefore a controlled approximation tied to a
separation of time scales.

\section{What Must Be Remembered}
\label{sec:hj-summary}

\begin{itemize}
\item The Hamilton-Jacobi equation is the condition that a canonical
transformation make the new Hamiltonian trivial.
\item A complete integral \(S(q,\alpha,t)\) contains the full solution, but
only when the mixed Hessian in \((q,\alpha)\) is nondegenerate.
\item In autonomous systems, \(S=W-Et\), and \(W\) solves
\(H(q,\partial W/\partial q)=E\).
\item Action variables are period integrals of the canonical one-form.
\item Adiabatic invariance is the slow-parameter version of action
conservation.
\end{itemize}

The conceptual bridge to the earlier chapters is this: Liouville integrability
produces invariant tori, action variables measure cycles on those tori, and
Hamilton-Jacobi theory constructs the canonical transformation that reveals
those variables when the system is separable enough.
